\section{Transference for dilation graphs}\label{sec:transference}

We transfer a norm estimate for independent prime labels to an averaged
moment of a dilation graph whose labels must divide the integer at their
vertex.  Averaging a required divisibility condition supplies a factor
$1/p$, which explains the reciprocal-prime law in the independent model.
A prime used at several vertices pays this cost only once.  The main
problem is to retain that dependence while using the independent-model
norm on most edges.  The resulting moment will control pairings in which
one endpoint is constant on all mark lists at its integer position.
Signed closed walks and repeated prime labels also occur in the
divisibility-graph arguments of
Helfgott and Radziwi{\l}{\l}~\cite[Section~1.5]{HelfgottRadziwill2021}
and Pilatte~\cite[Sections~3 and~5]{Pilatte2026}.
The operators and ranges here differ; the memory identity and
transference estimate below are proved for the present model.

\subsection{Prime groups and the two operators}

\begin{definition}[Prime groups]\label{def:prime-groups}
Fix constants
\[
 0<a<b<c<d<0.47,\qquad c_0,C_0,v_-,v_+>0,
 \qquad 0<q<1,\qquad \ell\in\N.
\]
For each sufficiently large $x$, let the finite, pairwise disjoint sets
of primes $\mathcal P_g$ be indexed by the disjoint sets
$\mathcal S$ and $\mathcal B$, where
\[
 c_0T\le |\mathcal S|,|\mathcal B|\le C_0T.
\]
The small and big groups, respectively, satisfy
\begin{align}
 \exp(L^a)\le p\le\exp(L^b)&\quad(p\in\mathcal P_g,
                 \ g\in\mathcal S),\notag\\
 \exp(L^c)\le p\le\exp(L^d)&\quad(p\in\mathcal P_g,
                 \ g\in\mathcal B),\label{eq:prime-group-ranges}\\
 v_-\le V_g:=\sum_{p\in\mathcal P_g}\frac1p&\le v_+,
 \qquad \mu_g(p):=\frac1{V_gp}.\notag
\end{align}
Write $s=|\mathcal S|+|\mathcal B|$ and
\[
 \omega_g(n)=\sum_{p\in\mathcal P_g}\ind_{p\mid n},
 \qquad \omega(n)=\sum_g\omega_g(n)
 \quad(n\in\Z).
\]
Thus divisibility at $n=0$ has its usual meaning.  We call primes in
$\bigcup_g\mathcal P_g$ group primes.
\end{definition}

Let $J$ be a sufficiently large fixed integer and put $M=J+\ell$.
An ordered label list has $M$ slots in each group.  A pattern $\nu$
specifies sets $C_g\subseteq\{1,\ldots,J\}$, with
\[
 C_g=\varnothing\quad(g\in\mathcal S),\qquad |C_g|\le m,
 \qquad t_g=\ell+|C_g|.
\]
The slots in $\{1,\ldots,J\}\setminus C_g$ are shared; their labels
are copied from the source list to the target list.  The other $t_g$
slots in each endpoint list are unshared.  To form $D$, use the shared labels
and, in the $C_g$ slots, independent labels of law $\mu_g$.  Thus
$D$ has $J$ prime factors from each group, counted with multiplicity.
Write $D=D_{\mathcal S}D_{\mathcal B}$.

For each pattern there is a complex coefficient $K_\nu$, depending
only on the ordered big-group source labels, target labels, and the
auxiliary free labels used in $D_{\mathcal B}$.  The finite pattern
family may depend on $x$ and $J$, but
\begin{equation}\label{eq:pattern-mass}
 \sum_\nu\sup|K_\nu|\le L^{C_1},
\end{equation}
where $m,C_1$ are fixed independently of $J$.  In every operator
below, symmetrization means averaging all $M!$ slot permutations
independently in each relevant group.  This is an orthogonal
projection because the measures are invariant under these permutations.
Our convention is that a row operator integrates a function at the
target and returns a function at the source.

\begin{definition}[Ideal operator]\label{def:ideal-operator}
On the space
\[
 \mathcal H_{\rm id}
 =L^2\left(\prod_{g\in\mathcal B}\mathcal P_g^M,
               \ \bigotimes_{g\in\mathcal B}\mu_g^{\otimes M}\right),
\]
repetitions of prime values are allowed.  For a real $\zeta$, the
unsymmetrized ideal row operation for $\nu$ retains shared labels,
samples all unshared target labels independently with laws $\mu_g$,
and samples the auxiliary labels of $D_{\mathcal B}$ independently
with the same laws.  Its multiplier is
\[
 K_\nu D_{\mathcal B}^{i\zeta}\ee{\Theta D_{\mathcal B}}.
\]
The sum over $\nu$, composed on both sides with big-group
symmetrization, is denoted $\mathcal T(\Theta)$.
\end{definition}

\begin{definition}[Physical operator]\label{def:physical-operator}
A physical state $(n,\mathbf p)$ consists of $n\in\Z$ and an
ordered list $\mathbf p=(p_{g,i})_{g,1\le i\le M}$ such that,
within each group, the labels are distinct and divide $n$.
Give every state at $n$ mass $\prod_gV_g^{-M}$ and use counting
measure in $n$.  This defines $\mathcal H_{\rm ph}$.

Fix an integer $k$ with $0<|k|\le L^{C_2}$.  For pattern $\nu$,
sample the auxiliary labels forming $D$ as above and move from
$n$ to $n'=n+kD$.  Sum over physical target states at $n'$ with
the prescribed shared labels, with coefficient
$\prod_g V_g^{-t_g}$ for each target.  Forbid every auxiliary free
label of $D$ from both endpoint lists.  Auxiliary free labels may
equal one another.  Multiply this row action by
\begin{equation}\label{eq:physical-multiplier}
 K_\nu D^{i\zeta}
 q^{(\omega(n)-Ms)/2}q^{(\omega(n')-Ms)/2}.
\end{equation}
The sum over patterns, symmetrized in all groups on both sides,
is $\mathcal A$.
\end{definition}

For clarity, the adjoint moves by $-kD$, interchanges source and
target labels in the coefficient, and conjugates the multiplier.
Indeed, after both lists and the auxiliary labels have been fixed,
the shared product and hence $D$ are unchanged on reversal.  The
joint measure of the two lists in group $g$ has normalization
\begin{equation}\label{eq:physical-joint-measure}
 V_g^{-M-t_g}
\end{equation}
in either direction.  The auxiliary-label probability is also
unchanged.  This proves the assertion before symmetrization, and
the symmetrizing projections are self-adjoint.

\begin{lemma}[Absolute physical bounds]\label{lem:physical-schur}
Replace the coefficients of each row transition of $\mathcal A$
by their absolute values before adding patterns or averaging
permutations.  The resulting operator and its adjoint have row
sums at most $L^{C_{\rm abs}}$, where $C_{\rm abs}$ may depend on
$m,C_1$ and the fixed data in \Cref{def:prime-groups}, but is
independent of $J$.  The same assertion holds if the position is
replaced by independent uniform residues at all group primes.
\end{lemma}

\begin{proof}
If the target has $y\ge M$ divisors from group $g$, then, after
the $M-t_g$ distinct shared labels have been fixed, there are
$(y-M+t_g)_{t_g}$ possible ordered unshared target lists.  Hence
their sum, with their endpoint damping, is at most
\begin{equation}\label{eq:physical-count-bound}
 \sup_{z\ge0}V_g^{-t_g}(z+t_g)_{t_g}q^{z/2}\le C.
\end{equation}
Here $t_g\le\ell+m$, so $C$ is independent of $J$.  The source
damping is at most one.  Free-label probabilities have total mass
one, and restrictions can only decrease an absolute sum.
Taking the product over $s=O(T)$ groups and using
\Cref{eq:pattern-mass} gives the row bound.  The reversed
normalization in \Cref{eq:physical-joint-measure} gives the identical
argument for columns.  Neither argument used any property of
integer positions beyond divisibility and translation.
\end{proof}

In particular all these operators are bounded on their stated Hilbert
spaces, including the physical space with unrestricted integer position.

\begin{theorem}[Local transference]\label{thm:transference}
For every fixed $E>0$ there is $E_{\rm id}>0$, depending only on
$E$ and the fixed data in \Cref{def:prime-groups}, with the following
property.  Assume
\begin{equation}\label{eq:ideal-norm-hypothesis}
 \sup_{\Theta\in\R}\norm{\mathcal T(\Theta)}\le L^{-E_{\rm id}}.
\end{equation}
Then, for every sufficiently large fixed $J$, where the lower
bound on $J$ may also depend on $m,C_1$, put
\[
 R=\lceil L^{0.52}\rceil,\qquad N=2R.
\]
For every fixed $\gamma>0$ and $I=[X,2X)$ with
$x^\gamma\le X\le x^{1/\gamma}$, one has
\begin{equation}\label{eq:transference-moment}
 \frac1{|I\cap\Z|}\sum_{n\in I\cap\Z}
 \langle u_n,(\mathcal A\mathcal A^*)^Ru_n\rangle
 \le L^{-EN},
\end{equation}
for all sufficiently large $x$.  Here $u_n$ is one on every
physical state at position $n$ and zero elsewhere.  The threshold
for $x$ may depend on all fixed parameters, including $J,C_2,
\gamma$, and the bound is uniform subject to these parameters
and \Cref{eq:ideal-norm-hypothesis}.  No bound on $\zeta$ is needed
apart from the assumed ideal estimate.
\end{theorem}

The vector $u_n$ is the constant function on the mark lists at $n$; it is
not normalized.  Thus the moment in \Cref{eq:transference-moment} sums
paths whose final integer position returns to $n$, with arbitrary initial
and final mark lists.  \Cref{cor:transference-pairing} will turn precisely
this estimate into cancellation against a mark-independent first
endpoint.  Those are the endpoints supplied in \Cref{sec:shifts}.

Here is the structure of the proof.  First replace the integer root by
independent residues and expand each big prime's contribution over
intervals containing all its active visits.  A memory records
its congruence between active uses.  The small-prime residues remain
physical: the many choices of omitted small marks make transfers between
memory and active lists rare.  On an edge with no such transfer, Fourier
transformation of the residue coordinates leaves the ideal big-label
operator.  We first allow different lifespans to use the same prime,
and then exclude such coincidences by inclusion--exclusion organized
by the number of independent equalities.  Only the rare-transfer
estimate needs the complexity-dependent choice of $J$; the ideal accuracy
is fixed beforehand.

\subsection{Replacing the root by independent residues}

Expand the left side of \Cref{eq:transference-moment} into paths
with $N$ alternating forward and backward edges.  Fix the
patterns, all endpoint permutations, the auxiliary labels, and
the active lists at visits $0,\ldots,N$.  These choices fix
offsets $h_0=0,h_1,\ldots,h_N$, with
\[
 h_{i+1}-h_i=\sigma_i kD_i,\qquad \sigma_i\in\{1,-1\},
 \qquad h_N=0.
\]
Only return to the position is imposed; the last ordered list
need not be the first one.  Let $\Lambda_i$ be the set of active
labels at visit $i$.  Put $q_0=q_N=\sqrt q$ and $q_i=q$ for
$0<i<N$.  The dependence on the initial integer $n$ is precisely
the nonnegative factor
\begin{equation}\label{eq:path-residue-factor}
 \prod_{i=0}^N\left(
 \prod_{p\in\Lambda_i}\ind_{n+h_i\equiv0\pmod p}\right)
 q_i^{\#\{p\text{ a group prime}:p\notin\Lambda_i,
                                      \ n+h_i\equiv0\pmod p\}}.
\end{equation}
The label normalization is initially $\prod_gV_g^{-M}$ and at
each edge $\prod_gV_g^{-t_g}$, in addition to the free-label
probabilities.  All remaining factors are independent of $n$.

\begin{lemma}[Uniform residue replacement]\label{lem:transference-crt}
For each compatible fixed path and every fixed $A>0$, the average
of \Cref{eq:path-residue-factor} over $I\cap\Z$ equals its
expectation under independent uniform residues at all group
primes, multiplied by $1+O(\exp(-ANT))$.  Incompatible active
congruences give zero in both models.
\end{lemma}

\begin{proof}
Let $Q_{\rm act}$ be the product of the distinct active primes
over all visits.  There are at most $(N+1)Ms$ of them, so
\begin{equation}\label{eq:active-crt-size}
 \log Q_{\rm act}=O_J(NTL^d)=o(L).
\end{equation}
Compatibility fixes one congruence modulo $Q_{\rm act}$.  All
damping factors belonging to these primes are then deterministic;
factor them out in both models.  This step permits a relative
estimate even when these deterministic factors are very small.

For a remaining prime, collect repeated offsets into distinct
residues.  Its factor has the form
\[
 1-X_p(n),\qquad
 X_p(n)=\sum_r c_{p,r}\ind_{n\equiv r\pmod p},
 \quad 0\le c_{p,r}\le1,
\]
with at most $N+1$ residues and $0\le X_p\le1$.  In the
independent model,
\[
 \lambda:=\sum_p\E X_p\le(N+1)\sum_p\frac1p=O(NT),
 \qquad \prod_p(1-\E X_p)\ge\exp(-O(NT)),
\]
since $\max_p\E X_p=o(1)$.  Bonferroni inequalities for numbers
in $[0,1]$ bracket $\prod_p(1-X_p)$ by consecutive truncations
of its elementary-symmetric expansion.  At degree $h$, the
independent expectation of the difference is at most
$\lambda^h/h!$.  Choose $h$ of order $NT$, with its fixed
constant sufficiently large after $A$.  Stirling's lower bound
$h!\ge(h/e)^h$ makes the difference smaller than
$\exp(-(A+C)NT)$ for any required fixed $C$.

Each truncated term specifies one residue at each of at most
$h$ additional primes.  Its CRT count on $I$ differs from its
independent density by $O(|I\cap\Z|^{-1})$, also when the
active congruence is imposed.  The number of such terms is at
most
\[
 \exp\big(O(h(L^d+\log N+\log s))\big)=\exp(o(L)),
\]
and their moduli, including $Q_{\rm act}$, have product
$\exp(o(L))$.  Here $0.52+d<1$.  Thus the total counting
error is $\exp(-\gamma L+o(L))$.  By
\Cref{eq:active-crt-size}, this is negligible relative to
$Q_{\rm act}^{-1}\exp(-O(NT))$.  Restoring the factored
deterministic damping proves the stated relative estimate.
\end{proof}

This replacement may be summed over paths with absolute
coefficients.  Indeed retain the full vector of residues as a
position, with its Haar probability measure, and translate it
by $\sigma_i kD_i$ on an edge.  The expected total mass of all
initial lists is at most one: in group $g$ it is
\begin{equation}\label{eq:initial-factorial-mass}
 V_g^{-M}\sum_{\substack{p_1,\ldots,p_M\in\mathcal P_g\\
                         \text{distinct}}}\frac1{p_1\cdots p_M}
 \le1.
\end{equation}
By \Cref{lem:physical-schur}, the absolute mass of all length
$N$ paths, even without a return condition, is $L^{O(N)}$.
Taking $A$ sufficiently large in \Cref{lem:transference-crt}
makes the summed error smaller than $L^{-A' N}$ for any desired
fixed $A'$.  We may therefore use the independent residue model.
The return condition is now imposed by the exact identity
\begin{equation}\label{eq:return-fourier}
 \ind_{h_N=0}=\int_0^1\ee{\theta h_N}\,d\theta.
\end{equation}
It suffices to bound the unrestricted signed path expression
uniformly in $\theta$.  Small-prime residues will remain actual
Haar coordinates throughout the argument.

\subsection{The exact expansion for one big prime}

Put $\eta_i=1-q_i$, $P_*:=\exp(L^c)$, and, for each big prime,
\begin{equation}\label{eq:lifespan-baseline}
 b_p=1-\frac1p\sum_{i=0}^N\eta_i>0,
 \qquad \widetilde\mu_g(p)=\frac1{V_gpb_p}.
\end{equation}
The ratio $\widetilde\mu_g(p)/\mu_g(p)$ is
$1+O(N/P_*)$, uniformly in big groups.  We extract the common
baseline $\prod_{p\ \text{big}}b_p\le1$ from every path sum.

Suppose first that $p$ is active at some visits.  Incompatible
active offsets give zero.  Otherwise let their common residue
be $c_p$, let $a_p,z_p$ be the first and last active visits, and
put $H_i=\ind_{h_i\equiv c_p\pmod p}$.  Its residue expectation
is
\[
 \frac1p\prod_{i:\ p\notin\Lambda_i}(1-\eta_iH_i).
\]
The two exact telescoping identities
\begin{align}
 \prod_{i<a_p}(1-\eta_iH_i)
 &=1-\sum_{j<a_p}\eta_jH_j
                     \prod_{j<i<a_p}(1-\eta_iH_i),\label{eq:prefix-ghost}\\
 \prod_{i>z_p}(1-\eta_iH_i)
 &=1-\sum_{j>z_p}\eta_jH_j
                     \prod_{z_p<i<j}(1-\eta_iH_i)\notag
\end{align}
follow by subtracting consecutive partial products.  Thus, in
addition to its active uses, the prime either starts at $a_p$
or has an earlier ghost birth at a hit $j<a_p$, with coefficient
$-\eta_j$; it either ends at $z_p$ or has a later ghost
termination at a hit $j>z_p$, again with coefficient $-\eta_j$.
At strictly internal unmarked hits it receives $q_i$.

If $p$ is never active, expand its full product and group each
term of degree at least two by its earliest and latest indices.
This gives
\begin{equation}\label{eq:orphan-expansion}
 b_p+\frac1p\sum_{j<i}\eta_j\eta_i
 \ind_{h_j\equiv h_i\pmod p}
 \prod_{j<t<i}\left(1-\eta_t
                         \ind_{h_t\equiv h_j\pmod p}\right).
\end{equation}
The empty and singleton terms give $b_p$; the interior product
is the sum over every possible further chosen index.  Relative
to the baseline, the prime is either absent or has an orphan
interval from $j$ to $i$, $j<i$, of cost
$(pb_p)^{-1}(-\eta_j)(-\eta_i)$ and the indicated internal
damping.

Consequently each represented big prime has exactly one
\emph{lifespan}: an interval containing all its active visits,
possibly containing none.  Its endpoints and all its active
visits have equal offsets modulo $p$.  It pays $(pb_p)^{-1}$
once, the ghost endpoint coefficients if present, and $q_i$
at each strictly internal unmarked hit.  No offset outside
the lifespan occurs in its weight.

There are two useful consequences of physical compatibility.
For large $x$ every group prime exceeds $|k|$.  If a prime is
active at both ends of an edge, then it divides $kD$, hence
$D$.  It cannot be an auxiliary free label, by the ban, so
it must be a prescribed shared label.  Conversely, a prime at
an unshared entry or exit cannot divide $D$: it would either
duplicate a shared label in that state or violate the free-label
ban.  Thus the original mark normalization assigns $V_g^{-1}$
to each maximal run of activity, and no additional factor at
shared continuation.  In particular, a lifespan with $r$ active runs
requires total label weight
\[
 \frac{V_g^{-r}}{pb_p}.
\]
The memory construction must therefore charge the reciprocal prime once,
at the first appearance, and charge $V_g^{-1}$ once for each active run.
Later congruence tests must not introduce further reciprocal-prime costs.
These requirements guide the choice of measures and transfer coefficients
below; the memory identity will verify them.

\subsection{A Hilbert space which remembers lifespans}

Let $Y_{\mathcal S}=\prod_{p\ \text{small}}\Z/p\Z$, with Haar
probability measure.  A small state at $y\in Y_{\mathcal S}$ is
an ordered list of $M$ distinct zero coordinates in each small
group, with mass $\prod_{g\in\mathcal S}V_g^{-M}$.  Write
$\mathcal H_{\mathcal S}$ for the resulting $L^2$ space.
For a big group set
\[
 \mathcal Z_g=\{(p,r):p\in\mathcal P_g, r\in\Z/p\Z\},
 \qquad d\lambda_g(p,r)=\widetilde\mu_g(p)
                                 \,d\operatorname{count}(r).
\]
Define its memory space by
\begin{equation}\label{eq:memory-hilbert-space}
 \mathcal F_g=\bigoplus_{j\ge0}
 L^2_{\rm sym}\left(\mathcal Z_g^j,\frac1{j!}\lambda_g^{\otimes j}\right),
 \qquad
 \mathcal H=\mathcal H_{\mathcal S}\otimes
 L^2\left(\prod_{g\in\mathcal B}\mathcal P_g^M,
                   \bigotimes_{g\in\mathcal B}\widetilde\mu_g^{\otimes M}\right)
 \otimes\bigotimes_{g\in\mathcal B}\mathcal F_g.
\end{equation}
Here ``symmetric'' means invariant under permutations of the
$j$ particle indices, including when their numerical values
coincide.  The $j=0$ summand is $\C$.  A pending particle
$(p,r)$ records displacement from its required hit, modulo $p$.
The residue measure in $\lambda_g$ is \emph{counting measure}:
requiring $r=0$ selects one point and incurs no factor $1/p$.
The boundary vector $\mathbf 1_\varnothing$ is one when every
memory is empty and zero otherwise; it is independent of all
other coordinates.

We next specify row actions.  The factorial measure in
\Cref{eq:memory-hilbert-space} is used for inner products and
adjoints, not inserted anew in every row transition.  All
operations are first understood on finite memory truncations;
the absolute bounds below justify the unrestricted sums.

Separate the evolution into an operation $G_i$ at each visit $0\le i\le N$
and an operation $\mathcal E_i$ across the following edge for $i<N$.
The visit operation handles ghost endpoints and unmarked hits.
The edge operation moves residues and transfers particles between active
lists and memory: a store keeps a departing active label pending, and a
promotion returns a pending particle to an active slot.  Thus $G_i$ acts
after arrival and promotion at visit $i$, and before storage and departure.

To make both operations bounded, choose $\rho\in(0,1)$ with
$\rho^2>\sqrt q$.  At a strictly internal unmarked hit, split the
required damping as
\[
 q_i=\rho\,(q_i/\rho^2)\,\rho.
\]
The adjacent edges supply the two factors $\rho$, and the visit operation
supplies the middle factor.  The edge factors will control the number
of possible promotions, while $q_i/\rho^2<1$ damps the visit operation.
The ghost endpoint factor is similarly split as
$-\eta_i=(-\eta_i/\rho)\rho$; the birth normalization is specified below.
The definitions implement these local factorizations, including the
initial and terminal visits.

For a group memory $\mathbf z=((p_h,r_h))_{h=1}^j$, write
$Z(\mathbf z)=\{h:r_h=0\}$.  At visit $i$ the ghost operation
$G_i$ is the product over big groups of the following row
action, leaving every other coordinate fixed:
\begin{align}
 (G_{i,g}F)(\mathbf z)
 ={}&\sum_{A\subseteq Z(\mathbf z)}
 \left(-\frac{\eta_i}{\rho}\right)^{|A|}
 \left(\frac{q_i}{\rho^2}\right)^{|Z(\mathbf z)|-|A|}
 \sum_{b\ge0}\frac1{b!}
 \left(-\frac{\eta_iV_g}{\rho}\right)^b\notag\\[-2mm]
 &\hspace{12mm}\cdot
 \int_{\mathcal P_g^b}
 F\big(\mathbf z\setminus A,(p'_1,0),\ldots,(p'_b,0)\big)
                 \prod_{h=1}^b d\widetilde\mu_g(p'_h).
 \label{eq:ghost-row}
\end{align}
Thus selected old zeros terminate, continuing zeros are damped,
and a new unordered batch is born at zero.  Termination precedes
birth; an orphan cannot be born and terminate at the same visit.

The edge operation $\mathcal E_i$, $0\le i<N$, sums the given
patterns in orientation $\sigma_i$, with independent source
and target symmetrizations in all active lists.  Between these
permutations its row action is as follows.
\begin{enumerate}
\item Form $D$ from shared source labels and the independent free
draws.  Retain the ban of free labels from both endpoint active
lists.  Multiply by the oriented coefficient
$K_\nu D^{i\zeta}$, or its conjugate with roles interchanged,
and by $\ee{\theta\sigma_i kD}$.
\item Multiply by $\rho$ for every old pending zero, before
storing any active particle.  In each big group select any
subset of its $t_g$ unshared source slots to store, appending
their particles at residue zero; drop the other unshared source
particles.  Shared particles remain active.
\item Translate $y$ and all pending residues, including the
newly stored ones, by $\Delta=\sigma_i kD$.
Choose any subset of the $t_g$ unshared big target slots and
an injection of those ordered slots into distinct \emph{old}
pending particle indices of that group whose translated residue
is zero.  Promote the selected particles into those slots,
remove them from memory, and multiply by $V_g^{-1}$ per
promotion.  Newly stored particles cannot be promoted across
this edge.  Require $p\nmid D$ for every stored or promoted
particle.  Fill the other unshared target slots by independent
fresh integrations against $\widetilde\mu_g$.
\item Multiply by $\rho$ for every pending zero remaining after
promotions.  In the small groups use the physical target sum,
with normalization $V_g^{-\ell}$, and the two small endpoint
factors $q^{(\omega_g-M)/2}$.  These counts concern small states
at $y$ and $y+\Delta$.
\end{enumerate}
Selections of promotion slots and injections are summed without
a factorial divisor.  All choices of source or target slots
refer to their canonical order between the sampled endpoint
permutations.  No big endpoint damping is included apart from
the pending-particle factors specified above.  Auxiliary free
draws of $D$ are not particle births.

For the moment impose, in the complete path sum, the rule that
all prime values assigned at \emph{different births} are
distinct.  Births comprise initial active slots, fresh active
target slots, and ghost births.  This is a global rule, even
if two births occur far apart.  Apart from it, place no
distinctness restriction on big active lists or memories.

\begin{lemma}[Exact memory identity]\label{lem:memory-identity}
With the distinct-birth rule, the independent-residue path
expression with the phase in \Cref{eq:return-fourier} equals
\begin{equation}\label{eq:memory-product}
 \left(\prod_{p\ \text{big}}b_p\right)
 \langle\mathbf1_\varnothing,
       G_0\mathcal E_0G_1\cdots\mathcal E_{N-1}G_N
       \mathbf1_\varnothing\rangle_{\rm distinct}.
\end{equation}
The subscript means that the indicator is inserted in the
expanded path integral, not that each local operator separately
enforces it.
\end{lemma}

\begin{proof}
A particle born active begins with its required hit at that
visit.  When it leaves an active run, storing it at zero and
then translating records subsequent displacement from that
hit.  Promotion checks displacement zero.  A later store
resets the anchor at a congruent offset, so it changes no
congruence requirement.  A ghost birth sets the anchor at an
unmarked hit, and a ghost termination also checks zero.
Shared continuation automatically respects the congruence.

The $\rho$ factors cancel locally.  At an internal visit a
continuing pending zero pays
$\rho(q_i/\rho^2)\rho=q_i$.  A ghost birth pays
$(-\eta_i/\rho)\rho=-\eta_i$, and a ghost termination pays
$\rho(-\eta_i/\rho)=-\eta_i$.  An active particle pays none
of these factors: promotion occurs before output damping and
storage after input damping.  At visit zero memory starts
empty, and at visit $N$ it must end empty, so there is no
missing boundary factor.  The values $q_0=q_N=\sqrt q$ are
therefore treated exactly.

For a prime with $r$ active runs, a first active birth costs
$1/(V_gpb_p)$ and each of its $r-1$ later promotions costs
$1/V_g$.  A first ghost birth costs $1/(pb_p)$ and its $r$
promotions cost $V_g^{-r}$.  Both give
\[
 \frac{V_g^{-r}}{pb_p}.
\]
An orphan has $r=0$.  These are precisely the baseline-relative
costs in \Cref{eq:prefix-ghost,eq:orphan-expansion},
including all original mark normalizations.

Conversely, fix an original compatible path and a chosen
lifespan term for every represented big prime.  Its active
visits determine exactly when it is shared, stored, promoted,
or dropped; its ghost endpoints determine its birth and
termination.  Between these uses it must remain pending.
The exclusions $p\nmid D$ follow from unshared entries and
exits as proved above, and immediate promotion of a newly
stored particle is impossible for the same reason.  A numerical
prime occupies at most one active slot or pending particle,
because it has only one birth.  Distinct ghost births in an
unordered batch are counted once by its factorial divisor.
Thus these constructions are mutually inverse and preserve
all coefficients.  This proves \Cref{eq:memory-product}.
\end{proof}

\Cref{fig:prime-lifespan} illustrates an admissible lifespan with
two runs of activity and ghost endpoints.

\begin{figure}[htbp]
\centering
\begin{tikzpicture}[x=1.25cm,y=0.85cm,>=Stealth,
  every node/.style={font=\small}]
 \draw[->,gray] (-.3,0)--(9.4,0) node[right]{visit};
 \foreach \x in {0,...,9}{\draw[gray] (\x,-.07)--(\x,.07);}
 \draw[very thick,dashed,blue!65!black] (0,0)--(2,0);
 \draw[very thick,blue!65!black] (2,0)--(3,0);
 \draw[very thick,dashed,blue!65!black] (3,0)--(6,0);
 \draw[very thick,blue!65!black] (6,0)--(7,0);
 \draw[very thick,dashed,blue!65!black] (7,0)--(9,0);
 \foreach \x in {2,3,6,7}{\fill[blue!65!black] (\x,0) circle (2.6pt);}
 \foreach \x in {0,9}{\draw[draw=blue!65!black,fill=white,very thick]
    (\x,0) circle (2.6pt);}
 \node[above] at (0,.1){ghost birth at $j$};
 \node[above] at (2.5,.1){active};
 \node[above] at (6.5,.1){active};
 \node[above] at (9,.1){terminate at $i$};
 \node[below,align=center] at (0,-.15){$-\eta_j$\\$1/(pb_p)$ once};
 \node[below,align=center] at (2,-.15){promote\\$1/V_g$};
 \node[below,align=center] at (3,-.15){store\\$r=0$};
 \node[below,align=center] at (6,-.15){promote\\$1/V_g$};
 \node[below,align=center] at (7,-.15){store};
 \node[below] at (9,-.15){$-\eta_i$};
 \node[above,align=center] at (4.5,.45){pending: $r\gets r+\Delta$\\
                                  unmarked zero pays $q_t$};
\end{tikzpicture}
\caption{One possible lifespan.  Solid segments denote shared
activity; dashed segments denote memory.  Every marked use and
ghost endpoint has the same offset modulo $p$.  The reciprocal
prime factor is charged only at birth; $\eta_t=1-q_t$ at visit $t$.}
\label{fig:prime-lifespan}
\end{figure}

\subsection{Absolute bounds, including the adjoint measures}

The memory identity is exact with the global distinct-birth rule.
We now allow different births to have equal prime values, so that each
operation acts locally on $\mathcal H$.  We will restore the global rule
after estimating this enlarged evolution.  The absolute bounds proved
here justify the memory truncation and remain available when equality
constraints later modify individual operations.  Their column bounds
require the factorial adjoint measures explicitly.

For a positive row kernel $K$, if a positive function $w$ satisfies
$Kw\le Rw$ and $K^*w\le Cw$, then
\begin{equation}\label{eq:weighted-schur}
 \norm K_{2\to2}\le\sqrt{RC}.
\end{equation}
Indeed, Cauchy's inequality in each row bounds $|Kf|^2$ by
$(Kw)K(|f|^2/w)$; integration and the column inequality prove
\Cref{eq:weighted-schur}.  We apply this to absolute kernels,
with
\[
 w(\text{state})=v^{j_{\rm tot}},
 \qquad j_{\rm tot}=\sum_{g\in\mathcal B}j_g,
\]
where $v>1$ is a fixed constant chosen below.

Here is a direct verification of the measures in the ghost
adjoint.  In one group abbreviate
\[
 A_i=\frac{q_i}{\rho^2},\qquad
 D_i=-\frac{\eta_i}{\rho},\qquad C_i=V_gD_i,
 \qquad \nu=\widetilde\mu_g,
\]
and let $z(\mathbf u)$ count zero residues in a memory list.
For test functions $f,h$, the joint integral from
\Cref{eq:ghost-row} is
\begin{align}
 \langle f,G_{i,g}h\rangle
 ={}&\sum_{k,t,b\ge0}\frac1{k!t!b!}
 \int A_i^{z(\mathbf u)}D_i^tC_i^b
       \overline{f(\mathbf u,\boldsymbol\alpha^0)}
       h(\mathbf u,\boldsymbol\beta^0)\notag\\
 &\hspace{14mm}\cdot
 d\lambda_g^{\otimes k}(\mathbf u)
 d\nu^{\otimes t}(\boldsymbol\alpha)
 d\nu^{\otimes b}(\boldsymbol\beta).
 \label{eq:ghost-joint-integral}
\end{align}
Here $\boldsymbol\alpha^0$ means that all its residues are zero.
The identity follows by selecting $t$ old zero indices from
a list of size $k+t$:
\[
 \frac1{(k+t)!}\binom{k+t}{t}=\frac1{k!t!}.
\]
Swapping $(t,\boldsymbol\alpha)$ with
$(b,\boldsymbol\beta)$ and conjugating shows that the adjoint
has the same row formula with birth coefficient $D_i$ and
termination coefficient $C_i$.  In particular, reversal moves
the factor $V_g$ from births to terminations; it does not create
any factor $p$ or $1/p$.

For the absolute ghost kernel the weighted row ratio is exactly
\begin{equation}\label{eq:ghost-schur-ratios}
 \exp\bigl(|C_i|v\nu(\mathcal P_g)\bigr)
 \left(A_i+\frac{|D_i|}{v}\right)^{z(\mathbf z)},
\end{equation}
and its column ratio has $C_i,D_i$ interchanged.  Choose $v$
large enough that, for both possible $q_i$,
\begin{equation}\label{eq:memory-weight-choice}
 \frac{q_i}{\rho^2}+
 \frac{2\max(1,v_+)\eta_i}{\rho v}<1.
\end{equation}
This is possible since $q_i\le\sqrt q<\rho^2$.  The same
choice works when every original ghost birth receives an
extra factor two.  Since $\nu(\mathcal P_g)=1+o(1)$, the
exponential factor in \Cref{eq:ghost-schur-ratios} is bounded
by a fixed constant per group.  Thus for some fixed $C_v$,
\begin{equation}\label{eq:ghost-absolute-bound}
 |G_i|w\le L^{C_v}w,\qquad |G_i|^*w\le L^{C_v}w,
 \qquad \norm{G_i}\le L^{C_v}.
\end{equation}
The exponent $C_v$ is independent of $m,C_1,J$.  All these
assertions include the version with doubled ghost births.

We verify the edge adjoint just as explicitly.  Fix one big
group, a pattern, endpoint permutations, free labels, and
the chosen transfer subsets.  Suppose $a$ target slots are
promotions and $b$ source slots are stores, so $a,b\le t_g$.
Let the source memory have size $j$, and write $u$ for its
$j-a$ surviving particles.  Let $p_1,\ldots,p_a$ be the
promoted prime labels and $s_1,\ldots,s_b$ the stored labels.
For step $\Delta$, the endpoint memories are
\begin{equation}\label{eq:edge-joint-memory}
 X=u\sqcup((p_h,-\Delta\bmod p_h))_{h\le a},
 \qquad
 Y=T_\Delta u\sqcup((s_h,\Delta\bmod s_h))_{h\le b},
\end{equation}
where $T_\Delta$ translates every residue.  The joint measure
on these variables is
\begin{equation}\label{eq:edge-joint-measure}
 \frac{d\lambda_g^{\otimes(j-a)}(u)}{(j-a)!}
 \prod_{h\le a}d\widetilde\mu_g(p_h)
 \prod_{h\le b}d\widetilde\mu_g(s_h),
\end{equation}
along with one $\widetilde\mu_g$ integration for every shared
active label, dropped source label, and fresh target label.
The transfer coefficient is $V_g^{-a}$, and the damping is
$\rho^{z(X)+z(Y)}$.  The factorial accounting is
\begin{equation}\label{eq:edge-factorial-reversal}
 \frac{(j)_a}{j!}=\frac1{(j-a)!}
 =\frac{(j-a+b)_b}{(j-a+b)!}.
\end{equation}
The first equality selects old memory indices into the $a$
specified ordered target slots.  The second selects $b$
indices of the output memory into the original source slots
on reversal.  Each required residue in
\Cref{eq:edge-joint-memory} selects exactly one point of
counting measure.  Translation preserves that measure.

It follows that the reversed row operation retrieves the
original stores and stores the original promotions; the
coefficient $V_g^{-a}$ remains unchanged.  If it is expressed
using the forward convention that charges $V_g^{-b}$ for its
reversed promotions, its additional factor is $V_g^{b-a}$.
This is bounded uniformly for $a,b\le\ell+m$.  Fixing shared
and free $D$-labels fixes $\Delta$ independently of which
memory indices are retrieved, as required for this change
of variables.  In small groups Haar translation and the joint
normalization $V_g^{-M-\ell}$ are invariant on reversal.

The identities above continue to hold with numerical
coincidences.  For example, a memory multiset with multiplicities
$n_z$ has mass $\prod_z\lambda_g(z)^{n_z}/n_z!$.  Subset
deletions and ordered retrievals retain their index
multiplicities.  No passage from indices to distinct numerical
values was made in \Cref{eq:ghost-joint-integral} or
\Cref{eq:edge-factorial-reversal}.

\begin{lemma}[Absolute memory bounds]\label{lem:memory-absolute}
There is $C_a$, allowed to depend on $m,C_1$ but independent
of $J$, such that every absolute edge has weighted row and
column bounds $L^{C_a}$.  These bounds and
\Cref{eq:ghost-absolute-bound} persist under arbitrary
multipliers of modulus at most one on individual local
transition choices, and under memory-size projections.
\end{lemma}

\begin{proof}
For a row, fix a pattern, free labels, permutations, and
transfer-slot subsets of sizes $a_g,b_g\le t_g$.  If $Z'$
old pending particles have zero translated residue in a
group, their ordered promotions and the remaining output
damping cost at most
\begin{equation}\label{eq:promotion-count}
 (Z')_{a_g}\rho^{Z'-a_g}
 \le\sup_{z\ge a_g}(z)_{a_g}\rho^{z-a_g}<\infty.
\end{equation}
Damping of stored particles can be discarded for this upper
bound.  The Schur weight contributes $v^{b_g-a_g}$, promotion
coefficients contribute $V_g^{-a_g}$, and fresh integrations
have mass $\widetilde\mu_g(\mathcal P_g)^{t_g-a_g}\le2^{t_g}$
for large $x$.  There are at most $4^{t_g}$ choices of transfer
subsets.  Each cost is bounded by a constant per group,
independent of $J$.  Small transitions satisfy
\Cref{eq:physical-count-bound}.  Multiplying over groups,
averaging permutations and free labels, and summing
\Cref{eq:pattern-mass} proves the row bound.

For the column, use \Cref{eq:edge-joint-memory,eq:edge-factorial-reversal} in reverse.  Original stores
are now ordered retrievals.  Their count is controlled by the
original input $\rho$ factors after reversed translation, in
exactly the form in \Cref{eq:promotion-count}.  The remaining
factor $V_g^{b_g-a_g}$ is bounded per group.  This proves the
column bound.  Restrictions and local multipliers of modulus
at most one are dominated by these absolute kernels.
\end{proof}

We now restrict the total memory size to
\begin{equation}\label{eq:memory-truncation}
 B=\lceil L^2\rceil
\end{equation}
before and after every operation, writing the projections
implicitly.  At most $(\ell+m)|\mathcal B|N$ particles can
be supplied by stores.  Thus a path violating this restriction
has at least $B-O((\ell+m)sN)$ ghost births.  Multiply every
ghost birth by two in the absolute path sum and use the
preceding bounds.  The total omitted mass, even with the
distinct-birth indicator, is at most
\begin{equation}\label{eq:memory-tail}
 2^{-B+O((\ell+m)sN)}L^{O(N)}
 =\exp(-\Omega(L^2))=o(L^{-AN})
\end{equation}
for every fixed $A$.  The empty-memory boundary vector has
bounded norm: the small part has mass at most one by
\Cref{eq:initial-factorial-mass}, and the big part has mass
\[
 \prod_{g\in\mathcal B}
 \widetilde\mu_g(\mathcal P_g)^M
 =1+O_J(sN/P_*)=1+o(1).
\]
These estimates also prove absolute summability before
truncation, for instance by first restricting all batch sizes
and then applying monotone convergence to absolute kernels.

\subsection{Rare transfers and the role of the small groups}

It remains to obtain a small signed edge norm and then restore the
omitted distinct-birth rule.  Decompose each edge as
$\mathcal E_i=\mathcal E_{i,\rm clean}+\mathcal E_{i,\rm dirty}$:
a term is \emph{dirty} if it has at least one store or promotion, and
\emph{clean} otherwise.  Dirty terms connect memory to the active lists;
we suppress them by averaging the omitted small labels.  Clean terms
only translate memory residues and refresh unshared active labels; the
next subsection will reduce their norm to the ideal estimate.  Define
\begin{equation}\label{eq:small-omission-number}
 \mathcal D=\binom{M}{\ell}^{|\mathcal S|}.
\end{equation}

\begin{lemma}[Rare transfer]\label{lem:rare-transfer}
With the memory truncation in \Cref{eq:memory-truncation}, the
aggregate of all dirty terms of any edge has norm at most
\begin{equation}\label{eq:dirty-norm}
 2L^{C_a}\left(\frac B{\mathcal D}\right)^{1/2},
\end{equation}
after enlarging $C_a$ independently of $J$ if necessary.
\end{lemma}

\begin{proof}
First consider the positive sum of terms with a promotion.
Fix its source, pattern, big-group permutations and shared
choices, and free $D$-labels.  For some one of its at most
$B$ old pending particles $(p,r)$, promotion requires
\begin{equation}\label{eq:promotion-congruence}
 r+\sigma_i kD_{\mathcal B}D_{\mathcal S}\equiv0\pmod p,
 \qquad p\nmid D.
\end{equation}
In each small group the source symmetrization omits an
independent uniformly chosen $\ell$-subset of its $M$
distinct labels.  If $F$ is the full product of all small
source labels and $U$ the product of the omitted labels,
then $D_{\mathcal S}=F/U$.  Unique factorization and
disjointness of groups show that the $\mathcal D$ choices
give distinct positive integers $U$.  Moreover,
\[
 U\le\exp(\ell|\mathcal S|L^b)<\exp(L^c)\le p
\]
for sufficiently large $x$.  They are therefore distinct
modulo $p$.  The quantities $F,k,D_{\mathcal B}$ are
invertible modulo every $p$ eligible in
\Cref{eq:promotion-congruence}.  Hence each pending particle
allows at most one omission choice.  At most $B$ of the
$\mathcal D$ equally likely choices allow any promotion.

Conditional on every omission choice, the remaining weighted
row costs are uniformly bounded by the proof of
\Cref{lem:memory-absolute}.  Thus the promotion part improves
its row bound to $L^{C_a}B/\mathcal D$, while its column
bound remains $L^{C_a}$.  Schur's test gives the first half
of \Cref{eq:dirty-norm}.  Terms with stores and no promotion
have the same improvement in their column bound by reversing
the edge: a stored particle becomes a retrieval from output
memory.  The exclusion $p\nmid D$ is retained under reversal,
even when numerical coincidences have been allowed.  Add
these two bounds.
\end{proof}

The coefficient restriction to big labels is used here:
conditioning on all such labels does not bias the small
omission choices.  In particular,
\begin{equation}\label{eq:dirty-J-gain}
 2L^{C_a}\sqrt{B/\mathcal D}
 \ll L^{C_a+1-(c_0/2)\log\binom{J+\ell}{\ell}}.
\end{equation}
An arbitrarily strong fixed saving is obtained by increasing
$J$, without changing the ideal norm target.

\subsection{Clean edges and Fourier reduction}

\begin{lemma}[Clean norm]\label{lem:clean-norm}
There is $C_s$, depending only on the fixed data in
\Cref{def:prime-groups}, such that the clean part of every
edge satisfies
\begin{equation}\label{eq:clean-norm}
 \norm{\mathcal E_{i,\rm clean}}
 \le 2L^{-E_{\rm id}+C_s}
\end{equation}
for all sufficiently large $x$, under
\Cref{eq:ideal-norm-hypothesis}.  In particular $C_s$ is
independent of $m,C_1,J$.
\end{lemma}

\begin{proof}
A clean edge leaves all memory prime lists and sizes
unchanged, translates their residues, drops every unshared
big source slot, and fills every unshared big target slot
freshly.  Its input and output $\rho$ factors and its
memory projections are contractions.  They commute with
active-list symmetrizations.  Absorb these factors and the
symmetrizations into the two test vectors $w_1,w_2$.
These vectors remain symmetric in the big active lists.

Fix an ordered tuple $\mathbf d$ of $J$ shared labels in
each small group and let $d'$ denote their joint product.
Define a map to unrestricted small residue functions by
\begin{equation}\label{eq:small-conditioning-map}
 (P_{\mathbf d}w)(y)=
 \left(\prod_{g\in\mathcal S}V_g^{-J/2-\ell}\right)
 q^{\frac12\sum_{g\in\mathcal S}(\omega_g(y)-M)}
 \sum_{\mathbf R:\ (\mathbf d,\mathbf R)\text{ a small state at }y}
             w(y,(\mathbf d,\mathbf R)).
\end{equation}
All other coordinates are left untouched.  The map is zero
when the displayed list is not a small state, so the damping
is used only where $\omega_g(y)\ge M$.  The two endpoint
factors in \Cref{eq:small-conditioning-map} give
$V_g^{-J-2\ell}=V_g^{-M-\ell}$, exactly the physical joint
normalization.

By Cauchy's inequality in $\mathbf R$, followed by summation
over $\mathbf d$, one has
\begin{equation}\label{eq:small-map-norm}
 \sum_{\mathbf d}\norm{P_{\mathbf d}w}^2
 \le L^{C_s}\norm w^2.
\end{equation}
To see the independence of $J$, the multiplier relative to
the original small state measure in group $g$ is bounded by
\[
 \sup_{z\ge M}
 q^{z-M}\frac{(z-M+\ell)_\ell}{V_g^\ell},
\]
a constant depending only on $q,\ell,v_-$.  Taking products
over $O(T)$ groups proves \Cref{eq:small-map-norm}.

The clean bilinear form is now exactly
\begin{equation}\label{eq:clean-conditioned-form}
 \sum_{\mathbf d}
 \langle P_{\mathbf d}w_1,
                \mathcal T'_{\mathbf d}P_{\mathbf d}w_2\rangle.
\end{equation}
The operator between these maps uses the big-list transitions
and translates all unrestricted small and memory residues
by $\sigma_i kd'D_{\mathcal B}$.  Fix the memory sizes and
their prime lists.  Fourier transformation on the finite
abelian group consisting of the small residue torus and
the pending residue coordinates diagonalizes these
translations.  Counting measure on pending residues and
Haar measure on the small torus each have the usual unitary
finite Fourier transform.  One can first use ordered memory
representatives; restriction to symmetric functions preserves
the bound.  Big active-list symmetry is unaffected.

At any fixed Fourier frequency, translation and the closure
factor contribute
\[
 \ee{\sigma_i\Theta D_{\mathcal B}}
\]
for a real $\Theta$ depending on $\mathbf d$, the memory
prime list, the frequency, and $\theta$, but independent
of the variable big active labels.  The scalar factor
$(d')^{\sigma_i i\zeta}$ has modulus one.  Consequently
the remaining operator, on symmetric big active tests, is
$\mathcal T(\Theta)$ or its adjoint, with two changes only:
the measures of active labels and fresh target integrations
are $\widetilde\mu_g$ in place of $\mu_g$, and the
free/active collision ban remains in force.

Both changes are negligible in operator norm.  Across the
$M|\mathcal B|$ active coordinates, the product density
changes by
\[
 1+O_J(sN/P_*).
\]
The corresponding square-root density identification is
unitary between the two $L^2$ spaces and commutes with
symmetrization.  Fresh integration densities obey the
same bound.  The ideal absolute row and column sums are
at most $L^{C_1}$ before these changes, since all its
unrestricted fresh laws are probabilities.  Schur's test
therefore bounds the density perturbation by
$O_J(L^{C_1}sN/P_*)$.

For the ban, in either row direction condition on the
fixed active state.  Each free auxiliary draw has maximal
atom $O(1/P_*)$; its chance to match a fixed active label
is bounded by that quantity.  A match to a fresh opposite
unshared label has the same bound by independence.
There are $O_{J,m}(s)$ possible comparisons.  Reversal
has the same estimate.  Thus removing the banned
transitions costs $O_{J,m}(L^{C_1}s/P_*)$ in norm.
These estimates are $o(L^{-A})$ for every fixed $A$.
They do not require exclusion of coincidences with
pending labels: such labels are fixed in the Fourier
decomposition and are already accounted for by translation.

It follows uniformly in $\mathbf d$ that the between-map
norm is at most $2L^{-E_{\rm id}}$ for large $x$.  Apply
Cauchy's inequality to \Cref{eq:clean-conditioned-form}
and then \Cref{eq:small-map-norm}; this proves
\Cref{eq:clean-norm}.
\end{proof}

Choose, in this order,
\begin{equation}\label{eq:transference-parameter-choice}
 G>E+C_v+3,\qquad E_{\rm id}>G+C_s+2,
\end{equation}
and finally choose the fixed integer $J$ so large that
\Cref{eq:dirty-J-gain} is smaller than
$\tfrac12L^{-G}$ for large $x$.  For example it suffices
that
\[
 \frac{c_0}{2}\log\binom{J+\ell}{\ell}>C_a+G+3.
\]
The clean estimate supplies the other half.  Hence the
whole signed truncated edge, with birth distinctness
still omitted, satisfies
\begin{equation}\label{eq:whole-edge-contraction}
 \norm{\mathcal E_i}\le L^{-G}.
\end{equation}
Only this final choice of $J$ needs to depend on $m,C_1$.
This establishes the required quantifier order, but a
further argument is necessary to impose distinct births
without losing the signed contraction on every edge.

\subsection{Restoring distinct births by equality rank}

A single absolute collision estimate cannot finish the proof: one
reciprocal-prime saving $P_*^{-1}$ does not absorb a length-$N$ absolute
path cost $L^{O(N)}$.  Instead, expand the global distinctness condition
by equalities between birth values.  Many independent equalities yield
many point-mass savings.  A small number of independent equalities can
be imposed by phases at their birth operations, preserving the signed
contraction on all the other edges.

In the truncated evolution allocate potential birth
addresses as follows: initial big slots; every possible
big target slot at every edge, in its canonical order
before final symmetrization; and indices $1,\ldots,B$
in each group's ordered ghost batch at each visit.
A target slot is performed as an address precisely when
it is a fresh birth rather than shared or promoted.
A ghost address is performed precisely when its batch
is at least that large.  These are conditions on the
choices made in their single local operation.  The
number $A_*$ of potential addresses satisfies
\begin{equation}\label{eq:birth-address-count}
 A_*\le M|\mathcal B|(N+1)+(N+1)|\mathcal B|B=L^{O(1)}.
\end{equation}
The implied exponent can be taken independent of $J$
once $x$ is sufficiently large for fixed $J$.

We use the following elementary form of distinctness
inclusion-exclusion.  Sum over collections $\pi$ of
disjoint blocks of potential addresses, every block
having size at least two.  Let $\mathcal C_\pi$ require
that all addresses in those blocks be performed and
that the prime values within each block be equal.  Set
\begin{equation}\label{eq:distinct-birth-expansion}
 \mu(\pi)=\prod_{B'\in\pi}(-1)^{|B'|-1}(|B'|-1)!,
 \qquad r(\pi)=\sum_{B'\in\pi}(|B'|-1).
\end{equation}
For a fixed realized path,
\begin{equation}\label{eq:distinct-indicator}
 \ind_{\text{distinct performed birth values}}
       =\sum_\pi\mu(\pi)\ind_{\mathcal C_\pi}.
\end{equation}
To verify this, interpret a block of size $h$ as a cycle
on its $h$ addresses: there are $(h-1)!$ such cycles,
each with sign $(-1)^{h-1}$.  The right side is the sum
of signs of permutations of the performed addresses
which preserve their numerical prime values.  It factors
over equal-value classes.  The sign sum is one for a
singleton class and zero for any class of size at least
two, by pairing permutations with their product by a
fixed transposition.

A rank-$r$ collection involves at most $2r$ addresses.
More precisely its total absolute coefficient, summed
over all collections of rank $r$, obeys
\begin{equation}\label{eq:rank-coefficient-count}
 \sum_{r(\pi)=r}|\mu(\pi)|\le A_*^{2r}\le L^{C'r}.
\end{equation}
For instance the unsigned cycle generating polynomial
on $A_*$ addresses is $\prod_{j=0}^{A_*-1}(1+jt)$,
where the power of $t$ is the sum of cycle lengths
minus the number of cycles.  Its coefficient of $t^r$
is at most $(\sum_jj)^r/r!\le A_*^{2r}$.

\paragraph{Large rank: chronological fresh integrations.}
Fix $\pi$.  Order birth addresses chronologically, using
canonical slot order inside an edge or batch, and call
the first address in each block its pivot.  Every other
address is a dependent birth.  When such an address is
performed, its fresh prime integration is constrained
to the already determined pivot value.  Since
\begin{equation}\label{eq:tilted-point-mass}
 \sup_{g\in\mathcal B,p\in\mathcal P_g}
                 \widetilde\mu_g(p)\ll P_*^{-1},
\end{equation}
this improves the weighted absolute row bound by
$O(P_*^{-1})$ per dependent birth.

Here the chronological assertion uses genuinely fresh
integrations.  At an edge choose stores, promotions,
and the step first, and then fill fresh target slots in
canonical order.  The step, promotion count, and memory
weight do not depend on these new prime values; the
coefficient is bounded by its supremum and exclusions
may be discarded.  In the proof of
\Cref{lem:memory-absolute}, replace the total fresh mass
in each constrained slot by its single-point bound.
For a ghost operation keep its ordered batch integral
and factorial divisor; constraining indicated fresh
coordinates replaces their mass factors by the same
single-point bound.  The birth coefficient $V_g$ and
the memory weight only give fixed constants per
constrained coordinate.  Initial active slots satisfy
the same product estimate.  Pivots and dependents
within one operation are handled in their prescribed
integration order.

For completeness, these conditional row bounds can be
iterated although a pivot may have died before its
dependent birth.  Retain all pivot values in the
conditioning.  Bound the remaining terminal sum by
backward induction using the uniform weighted row
bounds for every possible conditioning, beginning
with $\mathbf1_\varnothing\le w$.  Whenever a dependent
birth is reached, the preceding point-mass improvement
is uniform in its retained pivot value.  At the initial
boundary $w=1$, and the remaining integrated boundary
mass is bounded.  Thus the total absolute contribution
for this fixed collection is at most
\begin{equation}\label{eq:high-rank-path-bound}
 L^{C_h(N+1)}(C/P_*)^{r(\pi)},
\end{equation}
for fixed $C_h,C$; both may depend on $m,C_1$.

Choose a sufficiently large fixed $C''$, after these
constants and $C'$, and set
\begin{equation}\label{eq:rank-threshold}
 r_0=\left\lceil\frac{C''NT}{\log P_*}\right\rceil.
\end{equation}
By \Cref{eq:rank-coefficient-count} and
\Cref{eq:high-rank-path-bound}, the sum over $r\ge r_0$
is bounded by
\[
 L^{C_h(N+1)}\sum_{r\ge r_0}
               (CL^{C'}/P_*)^r=o(L^{-EN}).
\]
Indeed $\log P_*=L^c\gg T$, and increasing $C''$
beats the fixed exponent $C_h+E$.

\paragraph{Small rank: phases only at the affected births.}
For a block with pivot $\beta_0$ and other addresses
$\beta_1,\ldots,\beta_{h-1}$, equality of the integer
prime values has the exact representation
\begin{equation}\label{eq:birth-equality-fourier}
 \ind_{p_{\beta_j}=p_{\beta_0}\ (1\le j<h)}
 =\int_{[0,1]^{h-1}}
 \ee{\sum_{j=1}^{h-1}\tau_j(p_{\beta_j}-p_{\beta_0})}
             \,d\boldsymbol\tau,
\end{equation}
on paths performing all these addresses.  For fixed
Fourier parameters, collect the factors by their birth
operation.  Each affected operation is multiplied,
on its individual local choices, by
\[
 \ind_{\text{designated local births are performed}}
       \ee{\sum_{\beta\ \text{local}}c_\beta p_\beta},
\]
of modulus at most one.  Initial-slot factors merely
modify the initial boundary vector by such a multiplier.
An affected edge can lose its signed norm saving, but
\Cref{lem:memory-absolute} bounds its new norm by
$L^{C_a}$.  Unaffected edges retain their complete signed
pattern sum and both symmetrizations, so
\Cref{eq:whole-edge-contraction} still applies.

The ordered indices in a ghost batch do not require
ordered-memory functions.  For a batch of size $b$, insert
its local multiplier $\phi_b(p_1,\ldots,p_b)$ inside
the existing $1/b!$ product integral in
\Cref{eq:ghost-row}.  The rest of that integral is
symmetric in its fresh variables, so its value is
unchanged on replacing $\phi_b$ by
\[
 \frac1{b!}\sum_{\sigma\in S_b}
       \phi_b(p_{\sigma(1)},\ldots,p_{\sigma(b)}).
\]
This average has modulus at most one, also on numerical
diagonals.  The resulting operator acts on symmetric
memories, and its reversed multiplier is the conjugate
on the corresponding deleted batch in
\Cref{eq:ghost-joint-integral}.  Its absolute row and
column bounds remain \Cref{eq:ghost-absolute-bound}.
In particular no phase has to follow a particle through
later storage, propagation, or promotion: equality was
encoded entirely at its two birth addresses.

At most $2r$ edges are affected by a rank-$r$ collection.
The boundary norms are bounded, and all $N+1$ ghost
operations cost at most $L^{C_v}$.  Therefore its
Fourier-integrated contribution is at most
\begin{equation}\label{eq:low-rank-path-bound}
 O\left(L^{(-G+C_v)N+C_v+2r(G+C_a)}\right).
\end{equation}
Since
\[
 \frac{r_0}{N}=O(T/L^c)+O(N^{-1})=o(1),
\]
the sum over $r<r_0$, with the coefficient bound
\Cref{eq:rank-coefficient-count}, is
\[
 L^{(-G+C_v+o(1))N}=o(L^{-EN})
\]
by \Cref{eq:transference-parameter-choice}.

Thus the high-rank terms are controlled by their fresh-label costs,
while the low-rank terms retain signed contraction on $N-o(N)$ edges.
Together they bound the expansion with distinct births, which is the
expression required by \Cref{lem:memory-identity}.

\begin{proof}[Completion of \Cref{thm:transference}]
Apply \Cref{eq:distinct-indicator} in the truncated
memory identity.  The large- and small-rank estimates
bound it by $o(L^{-EN})$, uniformly in $\theta$.
Restore the tail \Cref{eq:memory-tail}; multiply by
the baseline, which is at most one; and integrate
\Cref{eq:return-fourier}.  Finally restore the
uniform residue-replacement error, choosing its
precision after the absolute path bound.  The result
is $o(L^{-EN})$, and therefore
\Cref{eq:transference-moment} for sufficiently large
$x$.  The choice of $E_{\rm id}$ in
\Cref{eq:transference-parameter-choice} depends only
on $E,C_v,C_s$, and the dependence of these constants
was stated above.  This completes the proof with
the asserted quantifiers.
\end{proof}

\subsection{The endpoint pairing consequence}

We now convert the root moment into the pairing estimate used in
\Cref{sec:shifts}.  Constancy on a mark fiber is needed only at the first
endpoint; the second endpoint can be any vector with the stated norm.

\begin{corollary}\label{cor:transference-pairing}
Under the hypotheses and choices of
\Cref{thm:transference}, let $f\in\mathcal H_{\rm ph}$
be supported on positions $I=[X,2X)$ and independent
of the chosen marks, with value $f_n$ at position $n$.
Suppose
\[
 \sup_{n\in I}|f_n|\le\exp(O(\sqrt L)),\qquad
 \norm f,\norm g\le X^{1/2}L^{C_3}.
\]
Then
\begin{equation}\label{eq:transference-pairing}
 |\langle f,\mathcal A g\rangle|
       \ll XL^{-E+2C_3}.
\end{equation}
Only the first endpoint must be independent of its
mark list.
\end{corollary}

\begin{proof}
Let $H\ge1$ bound every displacement $|kD|$.  Since
$D$ has $Js$ factors, all at most $\exp(L^d)$,
\begin{equation}\label{eq:maximum-graph-step}
 \log H=O_J(TL^d).
\end{equation}
Partition $I$ into consecutive intervals $I_j$ of
length $\lceil H\rceil$, except possibly the last.
Put $f_j=\ind_{I_j}f$ and restrict $g$ to the
$H$-neighborhood of $I_j$ to obtain $g_j$.  These
neighborhoods have bounded overlap, and
\[
 \langle f,\mathcal A g\rangle
       =\sum_j\langle f_j,\mathcal A g_j\rangle,
 \qquad \sum_j\norm{g_j}^2\ll\norm g^2.
\]
For $B_0=\mathcal A\mathcal A^*\ge0$, spectral
H\"older gives
\[
 |\langle f_j,\mathcal A g_j\rangle|
 \le\norm{g_j}\norm{f_j}^{1-1/R}
       \langle f_j,B_0^Rf_j\rangle^{1/(2R)}.
\]
The finite matrix
$(\langle u_n,B_0^Ru_m\rangle)_{n,m\in I_j\cap\Z}$
is positive semidefinite.  Its largest eigenvalue is
at most its trace, so mark independence implies
\[
 \langle f_j,B_0^Rf_j\rangle
 \le\left(\sum_{n\in I_j\cap\Z}|f_n|^2\right)d_j,
 \qquad
 d_j:=\sum_{n\in I_j\cap\Z}
                  \langle u_n,B_0^Ru_n\rangle.
\]
There is no assertion here that the $u_n$ have equal
norm, or that ordered lists return to their starting
values.  By \Cref{eq:maximum-graph-step},
\[
 \left(O(H)\sup|f_n|^2\right)^{1/(2R)}
 =\exp\left(O_J\left(\frac{TL^d+\sqrt L}{L^{0.52}}\right)\right)
 =O(1).
\]
Sum the resulting block bounds by H\"older with
exponents $2$, $2R/(R-1)$, and $2R$.  This gives
\[
 |\langle f,\mathcal A g\rangle|
 \ll\norm g\,\norm f^{1-1/R}
                       \left(\sum_jd_j\right)^{1/(2R)}.
\]
The theorem bounds the last sum by
$|I\cap\Z|L^{-2RE}\ll XL^{-2RE}$.
Substitution of the two endpoint norm bounds proves
\Cref{eq:transference-pairing}; the extra factor
$L^{-C_3/R}$ is bounded for fixed $C_3$.
\end{proof}
